\begin{recipe}[
    preparationtime = {10 min},
    portion = {\portion{2}},
    calory = {430kcal | 6g Protein | 34g Fett | 26g KH},
    price = {2,10€},
    rating = {1},
    created = {01.12.2025},
    source = {\protect\recipesource{https://www.instagram.com/p/DIjuqy2ODeX/}{Instagram: lass\_ma\_kochen}}
]{Acili Ezme}

\introduction{%
    Irgendwie hätte ich gedacht, dass es geiler wäre.
}


\ingredients{
    2 & Tomaten\index{Gemüse \& Pilze!Tomaten}\\
    1 & Zwiebel\index{Gemüse \& Pilze!Zwiebeln}\\
    2 & Rote Spitzpaprika\index{Gemüse \& Pilze!Spitzpaprika}\\
    2 & Grüne Spitzpaprika\index{Gemüse \& Pilze!Spitzpaprika}\\
    1--2 & Knoblauchzehen\index{Gemüse \& Pilze!Knoblauch}\\
    1 EL & Paprikamark\index{Vorrat \& Konserven!Paprikamark}\\
    1 EL & Tomatenmark\index{Vorrat \& Konserven!Tomatenmark}\\
    1 EL & Sumak\index{Kräuter \& Gewürze!Sumak}\\
    \nicefrac{1}{2} & Zitrone (Saft)\index{Obst!Zitronen}\\
    3 EL & Olivenöl\index{Öle, Saucen \& Würzmittel!Olivenöl}\\
    3 EL & Granatapfelsirup\index{Öle, Saucen \& Würzmittel!Granatapfelsirup}\\
    & Salz, Pfeffer\index{Kräuter \& Gewürze!Pfeffer}\index{Kräuter \& Gewürze!Salz}
}

\preparation{%
    \step%
    Alles grob zerkleinern und zusammen pürieren, bis eine dicke, aber nicht komplett flüssige Masse entsteht.

    \step%
    Die Masse durch ein Sieb oder mit einem Tuch gut ausdrücken, damit überschüssiger Saft entfernt wird.
    
    \step%
    Mindestens 1 Stunde im Kühlschrank ziehen lassen.
}

\suggestion[Serviervorschlag]{Mit Fladenbrot, als Grillbeilage oder zu Mezze servieren. Mit etwas zusätzlichem Olivenöl und Petersilie toppen.}

\end{recipe}